Il n'y a pas de réponse unique. Voici un exemple de travail attendu.

\emph{Énoncé proposé :} \og pour tout entier $n \geqslant 1$, $n^2 + n + 1$
est premier \fg.

\textbf{1.} $n = 1, 2, 3$ donnent $3$, $7$, $13$, qui sont premiers.

\textbf{2.} Mais $n = 4$ donne $21 = 3 \times 7$ : l'énoncé est faux.

\textbf{3.} Trois cas favorables ne couvrent que trois valeurs de $n$, alors
que l'énoncé en concerne une infinité ; seul un raisonnement valable pour
tout $n$ (ou un contre-exemple) permet de conclure.

On peut aussi proposer un énoncé vrai, par exemple \og la somme de trois
entiers consécutifs est divisible par $3$ \fg : le camarade doit alors le
démontrer, ici par $(n-1) + n + (n+1) = 3n$. \hfill $\blacksquare$
